\section{Research Interfaces and Deployment Modes}
\label{sec:interfaces-deployment}

GAWorld exposes the same persistent society through several interfaces, each
optimized for a different stage of research.  These interfaces are intentionally
thin over repository state and runtime commands: they improve inspection and
operation, but do not constitute independent simulation engines.

\subsection{Command Line and Interactive Inspection}

The command-line interface is the most direct research surface.  In addition to
\texttt{run}, \texttt{reset}, and \texttt{compare-event}, it can restrict the
simulation horizon, import external retrieval material, create an agent from
provided social-media text, serve visualization artifacts, and start local
services.  The \texttt{interview} command rebuilds a selected agent from the
configured population, restores its memory when statefulness is enabled, seeds
retrieval from that memory, and answers one or more researcher-supplied
questions.  Interviews are therefore probes of a stored agent state under a
new elicitation context; they are not independent measurements of the human
profile on which an agent may have been based.

LLM access is mediated by a provider router.  Current adapters cover local
Ollama, OpenAI-compatible, and Anthropic-style endpoints, while routing rules can
select providers by task or agent and can specify ordered fallbacks.  Calls are
logged with provider, task, agent, prompt size, latency, and fallback position.
This abstraction makes deployment adaptable to local or hosted models, but it
does not make providers behaviorally interchangeable.  A model or routing
change is an experimental change and should be recorded as such.

\subsection{Dashboard and Agent Studio}

The local Dashboard serves a static web client together with HTTP endpoints for
configuration, agent lists, profiles, state, memory, life events, simulation
start/stop status, trace metadata, and interviews.  It supports configuring
provider routing and initiating a run while exposing the resulting files to a
researcher.  Agent Studio extends this surface with an agent-centric view: it
can inspect and edit profiles and normalized state variables, show memory and
growth summaries, inspect skills and relationships, and launch an interview.
These tools make persistent state legible without requiring an analyst to read
every JSON or CSV file directly.

The web interface should not be conflated with the kernel's
\texttt{Controller}.  The latter validates structured movement and exposes
named, audited runtime interventions.  The Dashboard's primary run and mutation
paths still use HTTP handlers and CLI subprocess management rather than routing
all operations through that API.

\subsection{Trace Replay and External Environment Service}

The simulator's visualizer records a simulation trace and latest frame,
including agent positions, activities, events, and presentation assets.
\texttt{site/simviz} replays these artifacts in a browser and provides spatial
and temporal inspection distinct from aggregate plots.  Replay is valuable for
finding implausible schedules, stalled movement, or mismatches between an event
and an agent response.  It remains a view of what was recorded: omitted state
cannot be recovered from the animation, and visual plausibility is not a
validation criterion.

An optional external-environment server separates shared world evolution from
the main process.  Its HTTP API starts simulation days, advances time ticks,
returns environmental events and context, and persists engine state plus
day/tick caches.  It can use the configured LLM for environment generation or a
rule-based fallback.  This service boundary supports controlled reuse of a
world backend and makes environment state independently inspectable; the local
in-process environment remains a supported runtime path.

\subsection{Relay-Based Distributed Communication}

GAWorld also implements an opt-in distributed communication mode.  Simulation
nodes register local agents with an HTTP relay, discover agents in the same
cluster, poll bounded inbound messages, and send bounded outbound messages.
The relay persists its directory and retained message state, while the
simulation converts received messages into context for local cognition.  The
CLI can host the relay directly, and the Agent Studio can expose detected relay
activity.

This mechanism is best understood as a communication bridge, not as evidence
of scalable distributed simulation.  The current implementation does not
provide distributed clock synchronization, partitioned world-state ownership,
transactional consistency, fault-tolerant consensus, or published throughput
and scaling measurements.  It enables agents hosted by different processes or
machines to exchange messages through a shared service; claims about speedup,
population scale, or consistency under failure require separate evaluation.

\subsection{Deployment Boundaries}

The default services use Python threaded HTTP servers and file-backed state,
which keeps local research deployment transparent and easy to inspect.  Hosted
LLM credentials are supplied through environment variables rather than embedded
in profiles or manuscript artifacts.  Researchers exposing Dashboard,
environment, or relay services beyond a trusted machine must add production
authentication, transport security, access control, rate limiting, and data
governance appropriate to the profiles and generated content in use.  GAWorld's
current interfaces are research tooling; their existence should not be read as
a production-security or large-scale-operations claim.
